\section{Related Work and Research Gap}
\label{sec:related}

\paragraph{Literature search and positioning.} We positioned the work with a
structured search, not a systematic review. Twenty queries on SOC assessment, alert prioritization, security
process mining, peer analytics, supervisory technology and tamper-evident records
were run on OpenAlex and Crossref in September 2026 for publications from 2005
onwards (656 unique records, screened in a single pass; 106 abstracts
read; 51 included and joined with 26 sources of an earlier focused review), followed
by an update search for closest prior art and counterevidence. arXiv throttled the
automated queries and records beyond rank 25 per query were not screened, so absence
from this review is not evidence of absence. Queries, counts and decisions are in
the supplementary material.

\paragraph{SOC assessment and measurement.} Maturity models score capability through
self- or expert assessment~\cite{vanos2016soccmm,schlette2021ctisoc2m2}; metric
work defines and weights indicators of analyst or technical
performance~\cite{agyepong2019review,agyepong2023analyst,forsberg2023metrics}.
A case study in one academic SOC found that analysts who scored well on key
performance indicators did not necessarily further the SOC's goals~\cite{morin2025mismatched},
a caution for any metric-based supervision. SAIBERSOC measures an operational SOC
empirically by injecting synthetic attacks and observing outputs such as detection
accuracy and time to investigation, in an experiment with 124 students acting as
analysts~\cite{rosso2020saibersoc}. These approaches assess capability or test a SOC
actively; none analyses the records a SOC submits to an external supervisor.

\paragraph{Record-based assessment of incident handling (closest prior art).}
Palma et al.\ built a system that assesses an organization's incident-management
process for compliance with a reference process derived from ITIL and ISO/IEC~27035:
it aligns the event log with the model, classifies deviations with a taxonomy,
prioritizes them with a cost model and supports an auditor; it was validated on a
benchmark and demonstrated on a public real incident-management log~\cite{palma2024compliance}.
A companion benchmark compares such assessment models~\cite{palma2024benchimp}, and
two visual-analytics tools support the human assessor~\cite{palma2024eurova,palma2025impavid}.
The public log used in that line of work is the ServiceNow incident log (141,712 events
over 24,918 incidents) released with~\cite{amaral2019completion} and published
as~\cite{amaral2018uci498}---the same log we use for our external evaluation, and one
on which SLA attainment has also been modelled directly~\cite{swain2022sla}. Automated,
record-based assessment of incident handling for an assessor therefore already exists
and has been demonstrated on our external data. Table~\ref{tab:related} contrasts it
with SAT-SA.

\begin{table}[t]
\caption{Closest prior work compared with SAT-SA, as described in the cited sources
(``--'' not addressed; ``n/d'' not determinable from the source).}
\label{tab:related}
\centering
\scriptsize
\setlength{\tabcolsep}{2.5pt}
\begin{tabular}{@{}p{2.1cm}p{2.1cm}p{2.1cm}p{1.7cm}p{1.5cm}p{1.5cm}@{}}
\toprule
Work & Evidence & Analysis & Across organizations & Human role & Decision binding \\
\midrule
Maturity models \cite{vanos2016soccmm,schlette2021ctisoc2m2} & self / expert assessment & maturity scoring & comparable scores & assessor & -- \\
SAIBERSOC \cite{rosso2020saibersoc} & outputs of an operational SOC & synthetic attack injection & SOC configurations & analysts & -- \\
Palma et al.\ \cite{palma2024compliance,palma2025impavid} & incident-management event log & alignment, deviation cost model & n/d & auditor, assessor & -- \\
DomainPrio \cite{chiba2022domainprio} & one SOC's investigations & population analysis, regression & -- & analysts & -- \\
SAT-SA (this work) & six categories of submitted SOC records & rule-based workers, saturating fusion & peer cohorts, ranking & supervisor decides & signed receipt, ledger \\
\bottomrule
\end{tabular}
\end{table}


\paragraph{Alert and incident prioritization.} A large literature ranks alerts or
incidents for analysts~\cite{jalalvand2024prioritisation,tariq2025alertfatigue}:
resource-aware selection of alerts to investigate~\cite{shah2019alertselection},
provenance-based triage~\cite{hassan2019nodoze}, guided response learned from
customer triage labels~\cite{freitas2024guide}, learning to defer to
analysts~\cite{jalalvand2025l2d}, and, recently, LLM agents embedded in SOC
triage~\cite{banstola2026agentic}. DomainPrio analyses the population of a SOC's
investigations rather than single alerts~\cite{chiba2022domainprio}. SAT-SA differs in
unit and user---it ranks monitored organizations for a supervisor from findings about
how work was handled---and is not compared with these systems.

\paragraph{Process conformance and peer comparison.} Conformance checking relates
recorded behaviour to a model and exposes skipped or out-of-order
activities~\cite{rozinat2008conformance,carmona2018conformance,vanderaalst2016pm},
and has been applied to security audit trails~\cite{vanderaalst2005security} and to
IT incident management against ITIL~\cite{ferreira2008itil}. SAT-SA's execution-gap
and negative-space rules are rule-encoded conformance checks, not a new technique.
Peer group analysis flags entities that diverge from similar
entities~\cite{weston2008pga,bolton2002fraud}; SAT-SA uses median and median absolute
deviation (MAD)~\cite{leys2013mad} over a declared cohort.

\paragraph{Supervision and human oversight.} Financial supervisors use technology to
analyse submitted data (``suptech'') while keeping human judgement
central~\cite{broeders2018suptech}. Automation bias is a documented risk of decision
support~\cite{parasuraman2010automation} and has been reviewed for SOCs
specifically~\cite{tilbury2024automationbias}; we did not study it.

\paragraph{Integrity of records.} Tamper-evident logs detect alteration of recorded
history~\cite{schneier1999logs,crosby2009logging}; practical systems protect audit logs
with trusted execution~\cite{paccagnella2020custos} or ratcheting~\cite{guardiola2022sealfs},
and post-quantum primitives have been applied to forensic logging~\cite{fugkeaw2026pqabl}.
SAT-SA proposes no logging scheme or cryptography; it composes
SHA3-256~\cite{nist2015fips202} and ML-DSA-65~\cite{nist2024fips204} to bind one
canonical supervisory record to the authorized decision.

\paragraph{Gap.} Record-based incident assessment, supervisory analysis of submitted
data, peer comparison and tamper-evident records each exist. Within the databases,
queries, screening procedure and time period described above, we did not identify a
study that combines them for supervision of many SOCs from their evidence submissions---with
record-cited findings, cross-organization ranking and a cryptographically bound human
decision---and evaluates the combination against simple baselines and on independent
data. This is a search-bounded observation about a combination of established parts,
not a claim that no such study exists.
